\id{МРНТИ 65.59.19}{}

\begin{header}
\swa{Л.Б.~Умиралиева, М.А.~Дибирасулаев, Д.М.~Дибирасулаев, Л.Г.~Стоянова, C.С.~Турлыбаева, И.Д.~Филатов}{ИССЛЕДОВАНИЕ ВЛИЯНИЯ ПИЩЕВОГО ПОКРЫТИЯ НА ПОТЕРИ МАССЫ И ИЗМЕНЕНИЕ МИКРОБИОЛОГИЧЕСКИХ ПОКАЗАТЕЛЕЙ ТУШЕК ПТИЦЫ В ПРОЦЕССЕ ОХЛАЖДЕНИЯ И ХРАНЕНИЯ ПРИ БЛИЗКРИОСКОПИЧЕСКОЙ ТЕМПЕРАТУРЕ}

\tsp{1}Л.Б.~Умиралиева\envelope,
\tsp{2}М.А.~Дибирасулаев,
\tsp{2}Д.М.~Дибирасулаев,
\tsp{3}Л.Г.~Стоянова,
\tsp{1}C.С.~Турлыбаева,
\tsp{1}И.Д.~Филатов
\end{header}

\begin{affil}
\tsp{1}ТОО «Казахский научно-исследовательский институт перерабатывающей и пищевой промышленности» Алматы, Казахстан,

\tsp{2}Всероссийский научно-исследовательский институт холодильной промышленности - филиал ФГБНУ «ФНЦ пищевых систем им. В.М. Горбатова» РАН, Москва, Россия,

\tsp{3}МГУ им. М. В. Ломоносова, Москва, Россия

\corrauthor{Корреспондент-автор: l.umiraliyeva@rpf.kz}
\end{affil}

В данной работе представлены результаты исследований по обоснованию
оптимального состава пленкообразующих покрытий (ППС) для хранения
охлажденного мяса птицы на основе моноглицеридов и молочной сыворотки с
применением биоконсервантов. Основное внимание уделено оценке влияния
пищевых покрытий на потери массы (усушку) мяса птицы в процессе
охлаждения и последующего хранения при близкриоскопической температуре в
течение 30 суток. Производственная апробация технологии проводилась на
базе предприятия СПК «Медеу 55» (Туркестанская область). В ходе
испытаний применялся состав, включающий молочную сыворотку, глицерин,
органические кислоты (молочную и лимонную), а также бактериоцин
гибридного штамма \emph{Lactococcus lactis} ssp. \emph{lactis} F-116 в
качестве активного антимикробного компонента. Результаты испытаний
показали, что внедрение данной технологии позволяет не только повысить
микробиологическую безопасность продукции и продлить сроки её годности,
но и существенно снизить естественную убыль массы тушек птицы.
Дополнительно в течение 30 суток хранения изучалось влияние различных
комбинаций ППС на общее микробное число (ОМЧ). В частности, сравнивалась
эффективность применения состава на основе 0,1\% низина «Nisin Рro»
совместно с 1\% желатином и 2\% молочной кислотой и состава с
использованием 0,1\% культуральной жидкости штамма \emph{Lactococcus
lactis} ssp. \emph{lactis} F-116.

{\bfseries Ключевые слова:} мясо птицы, охлаждение, пищевые покрытия,
потери массы, близкриоскопическая температура, микробиологические
показатели, хранение мяса птицы

\begin{header}
INVESTIGATION OF THE EFFECT OF FOOD COATING ON WEIGHT LOSS AND CHANGES IN MICROBIOLOGICAL PARAMETERS OF POULTRY CARCASSES DURING COOLING AND STORAGE AT NEAR-CRYOSCOPIC TEMPERATURE

\tsp{1}L.B.~Umiraliyeva\envelope,
\tsp{2}M.A.~Dibirasulaev,
\tsp{2}D.M.~Dibirasulaev,
\tsp{3}L.G.~Stoyanova,
\tsp{1}S.S.~Turlybayeva,
\tsp{1}I.D.~Filatov
\end{header}

\begin{affil}
\tsp{1}Kazakh Scientific Research Institute of Processing and Food Industry LLP, Almaty Kazakhstan

\tsp{2}All-Russian Scientific Research Institute of Refrigeration Industry - branch of V.M. Gorbatov Federal Research Center for Food Systems of Russian Academy of Science, Moscow, Russia,

\tsp{3}Moscow State University named after M.V. Lomonosov, Moscow, Russia,

е-mail: l.umiraliyeva@rpf.kz
\end{affil}

This study presents research results on justifying the optimal
composition of edible film-forming coat\-ings (EFC) for the storage of
chilled poultry meat, based on monoglycerides and whey incorporated with
biopreservatives. The primary focus is placed on evaluating the effect
of edible coatings on the mass loss (shrinkage) of poultry during
cooling and subsequent storage at sub-cryoscopic temperatures for a
period of 30 days. Industrial validation of the technology was conducted
at the "Medeu 55" Agricultural Production Cooperative (Turkistan
region). The tested formulation included whey, glycerin, organic acids
(lactic and citric), and a bacteriocin from the hybrid strain
\emph{Lactococcus lactis} ssp. \emph{lactis} F-116 as an active
antimicrobial agent. The results demonstrated that the implementation of
this technology not only enhances the microbiological safety of the
product and extends its shelf life but also significantly reduces the
natural mass loss of poultry carcasses. Additionally, the impact of
various EFC combinations on the total viable count (TVC) was studied
over a 30-day storage period. Specifically, the efficacy of a 0.1\%
Nisin Pro-based composition combined with 1\% gelatin and 2\% lactic
acid was compared to a formulation utilizing 0.1\% culture liquid from
the \emph{Lactococcus lactis} ssp. \emph{lactis} F-116 strain.

{\bfseries Keywords:} poultry meat, cooling, edible coatings, weight loss,
near-cryoscopic temperature, microbio\-logical indicators, poultry storage

\begin{header}
ЖАҚЫН КРИОСКОПИЯЛЫҚ ТЕМПЕРАТУРАДА САЛҚЫНДАТУ ЖӘНЕ САҚТАУ ҮДЕРІСІНДЕ ҚҰС ҰШАЛАРЫНЫҢ МАССА ЖОҒАЛТУЫ МЕН МИКРОБИОЛОГИЯЛЫҚ КӨРСЕТКІШТЕРІНІҢ ӨЗГЕРУІНЕ ТАҒАМДЫҚ ҚАПТАМАНЫҢ ӘСЕРІН ЗЕРТТЕУ

\tsp{1}Л.Б.~Умиралиева\envelope,
\tsp{2}М.А.~Дибирасулаев,
\tsp{2}Д.М~Дибирасулаев,
\tsp{3}Л.Г.~Стоянова,
\tsp{1}C.С.~Турлыбаева,
\tsp{1}И.Д.~Филатов
\end{header}

\begin{affil}
\tsp{1}"Қазақ қайта өңдеу және тамақ өнеркәсібі ғылыми-зерттеу институты" ЖШС, Алматы Қазақстан,

\tsp{2}Тоңазытқыш Өнеркәсібінің Бүкілресейлік Ғылыми-Зерттеу Институты - В.М. Горбатов атындағы Ресей Ғылым Академиясының Азық-Түлік жүйелері Федералды ғылыми-зерттеу орталығының филиалы Мәскеу, Рессей,

\tsp{3}М. В. Ломоносов атындағы Мәскеу мемлекеттік университеті, Мәскеу, Рессей,

е-mail: l.umiraliyeva@rpf.kz
\end{affil}

Бұл жұмыста биоконсерванттарды қолдана отырып, моноглицеридтер мен сүт
сарысуы негізінде салқындатылған құс етін сақтауға арналған үлдір
түзетін жабындардың (ҮТЖ) оңтайлы құрамын негіздеу бойынша зерттеу
нәтижелері ұсынылған. Негізгі назар құс етін салқындату және одан әрі 30
тәулік бойы криоскопиялық температураға жақын жағдайда сақтау процесінде
тағамдық жабындардың салмақ жоғалтуына (кебуіне) әсерін бағалауға
аударылған. Технологияны өндірістік сынақтан өткізу «Медеу 55» АӨК
(Түркістан облысы) базасында жүргізілді. Сынақтар барысында құрамында
сүт сарысуы, глицерин, органикалық қышқылдар (сүт және лимон), сондай-ақ
белсенді микробқа қарсы компонент ретінде \emph{Lactococcus lactis}
ssp\emph{. lactis} F-116 гибридті штаммының бактериоцині бар композиция
қолданылды. Сынақ нәтижелері бұл технологияны енгізу өнімнің
микробиологиялық қауіпсіздігін арттыруға және оның жарамдылық мерзімін
ұзартуға ғана емес, сонымен қатар құс ұшаларының табиғи салмақ жоғалтуын
айтарлықтай төмендетуге мүмкіндік беретінін көрсетті. Қосымша 30 тәулік
сақтау кезеңінде ҮТЖ-ның түрлі комбинацияларының жалпы микробтық санына
(ЖМС) әсері зерттелді. Атап айтқанда, 1\% желатин және 2\% сүт
қышқылымен бірге 0,1\% «Nisin Рro» низині негізіндегі құрамның
тиімділігі \emph{Lactococcus lactis} ssp\emph{. lactis} F-116 штаммының
0,1\% культуралдық сұйықтығын қолданған құраммен салыстырылды.

{\bfseries Негізгі сөздер:} құс еті, салқындату, тағамдық қаптамалар, масса
жоғалту, жақын криоскопиялық температура, микробиологиялық көрсеткіштер,
құс сақтау.

\begin{multicols}{2}
{\bfseries Введение.} Анализ научных публикаций по технологии холодильного
консервирования мяса птицы показывает, что наряду с ин­тенсифика­цией
способов холодильной обработки, требующих значительных капитальных
затрат на техперевооружение камер холодильни­ков, одним из эффективных
путей снижения потерь массы и сохранения качества является обработка
поверхности продукта пище­выми пленкообра­зующими составами (ППС).

В основе массопереноса при охлаждении лежит механизм испарения влаги с
поверхности продукта в воздушную среду с её последующей конденсацией на
охлаждающих элементах оборудования. В случаях, когда испарение
происходит из сложных растворов или затруднено наличием полупроницаемых
барьеров (например, защитных пленок), традиционный коэффициент испарения
для чистой воды \emph{b\tsb{w}} заменяется на коэффициент
сопротивления испарению.

Следовательно, охлаждение пищевых продуктов в негерметичной упаковке
всегда сопровождается активным тепломассообменом, который влечет за
собой не только снижение температуры, но и неизбежную потерю массы.
Несмотря на это, именно охлаждение признано наиболее щадящим методом
консервирования. В сравнении с другими способами обработки, оно
оказывает минимальное воздействие на структуру продукта, позволяя
максимально долго сохранять его исходную свежесть и товарный вид.

Проблема естественной убыли массы из-за испарения влаги наиболее
критична для неупакованной продукции. Итогом такого обезвоживания
становится неизбежное снижение качественных характеристик сырья. В
качестве одного из наиболее эффективных решений выступает обработка мяса
и птицы антимикробными пленками. Обладая высокими барьерными свойствами,
такого вида покрытия минимизируют массообмен и удерживают влагу, что
позволяет существенно продлить сроки годности и сохранить товарный вид
продукта.

Анализ актуальных исследований в области мясной промышленности
подтверждает, что потери веса и ухудшение микробиологических показателей
остаются главными вызовами при хранении охлажденной птицы. Данные
негативные процессы спровоцированы интенсивной миграцией влаги,
деструкцией белковых соединений и активным ростом микрофлоры, включая
сапрофитные и условно-патогенные микроорганизмы {[}1-10{]}. Несмотря на
то, что использование температурного режима, близкого к точке замерзания
(от -1 до +1 °C), эффективно замедляет биохимическую порчу продукции,
применение подобных режимов требует внедрения дополнительных мер по
защите поверхности мяса птицы {[}2{]}.

Для повышения сохранности мяса перспективным направлением является
применение съедобных покрытий (edible coatings). Создаваемые на основе
природных полимеров - белков, полисахаридов и липидов - они образуют на
поверхности продукта физико-химический барьер, который не только
замедляет испарение влаги, но и препятствует окислению жиров, а также
ограничивает контакт продукта с внешней микробиотой {[}3{]}.
Установлено, что введение в состав пленок функциональных добавок
(антиоксидантов, антимикробных агентов, наночастиц) существенно
усиливает их защитные свойства, эффективно ингибируя размножение
микроорганизмов {[}3, 4{]}.

При хранении мяса птицы в условиях холодильника использование покрытий
на основе хитозана, альгината или казеина позволяет значительно снизить
потери влаги и подавить рост микроорганизмов по сравнению с
необработанными образцами, одновременно сохраняя привлекательные
органолептические свойства {[}4{]}. Аналогичный результат зафиксирован и
при включении в состав покрытий эфирных масел или природных экстрактов с
выраженной антибактериальной активностью {[}5{]}.

Несмотря на наличие широкой доказательной базы, в современной литературе
практически отсутствуют систематизированные данные о применении
антимикробных покрытий в условиях близкриоскопических температур. Между
тем, сочетание фактора глубокого охлаждения с барьерным действием
покрытий может дать значительный синергетический эффект, требующий
детального изучения.

Согласно результатам исследований {[}6-7{]}, синергия
низкотемпературного хранения и применения пищевых покрытий позволяет
эффективнее ингибировать окисление липидов и обеспечивать
микробиологическую стабильность продукции на более длительный срок, чем
использование одного лишь холода.

При изучении микрбиоты охлажденной птицы критическое значение имеет не
столько общая численность микроорганизмов, сколько качественный состав
доминирующих видов. Особую опасность представляют штаммы, способные
образовывать биопленки и сохранять активность при низких температурах
{[}7{]}. В совокупности возникает необходимость комплексного подхода к
мониторингу: важно отслеживать динамику развития конкретных таксонов,
непосредственно ответственных за порчу продукта, а не ограничиваться
только общими показателями.

Естественная убыль массы птицы продолжается и на этапе холодильного
хранения. В её основе лежат два ключевых фактора: поверхностное
испарение влаги и отделение тканевой жидкости. Темпы потери веса
напрямую зависят от длительности хранения и стабильности микроклимата в
камере {[}8-10{]}. Для минимизации усушки оптимальным признан режим с
температурой 0-2 °C и относительной влажностью воздуха в пределах 80--85
\%. Любое отклонение от данных параметров неизбежно провоцирует
ускоренную дегидратацию продукта.

Наличие и тип упаковки во многом определяет интенсивность усушки
продукции. При хранении птицы в открытом виде наблюдаются максимальные
потери веса из-за беспрепятственного испарения влаги. Применение
пленочных материалов, вакуумирования или модифицированной газовой среды
(МГС) позволяет не только существенно минимизировать потерю массы, но и
стабилизировать качественные характеристики птицы на протяжении всего
периода хранения {[}11-12{]}.

Сегодня в пищевой биотехнологии прослеживается отчетливый тренд на поиск
эффективных антибактериальных агентов, способных стать полноценной
альтернативой или дополнением к традиционным химическим консервантам и
антибиотикам. Особый интерес в данном контексте представляют
бактериоцины-пептидные соединения рибосомального синтеза, продуцируемые
различными группами бактерий.

Наиболее востребованным и коммерчески успешным представителем этой
группы является низин. Данный полициклический пептид, вырабатываемый
молочнокислыми бактериями \emph{Lactococcus lactis}, обладает выраженным
ингибирующим действием против широкого спектра грамположительных
микроорганизмов, включая опасные патогены и спорообразующие бактерии.
Благодаря своей доказанной безопасности, низин широко применяется в
качестве биоконсерванта во многих странах, включая США и страны Европы
{[}13-18{]}. Высокий потенциал данного бактериоцина открывает широкие
возможности для его интеграции в масштабные производственные линии и
дальнейшей диверсификации способов его применения.

Несмотря на большое количество преимуществ низина, его широкое внедрение
ограничивается рядом технологических недостатков природных
штаммов-продуцентов. Ключевыми проблемами остаются нестабильность
получаемого препарата, его высокая чувствительность к изменениям внешней
среды (в частности, к колебаниям pH и температуры), а также низкая
эффективность при попытках промышленного масштабирования биосинтеза.
Поскольку исходные штаммы синтезируют бактериоцины в малых
концентрациях, возникает острая необходимость в создании более
продуктивных биотехнологических решений. В этом плане определяющую роль
играют методы генетической инженерии и точная оптимизация условий
культивирования. Использование рекомбинантных штаммов на базе \emph{L.
lactis} позволяет преодолеть природные лимитирующие факторы и достичь
качественно новых показателей выхода целевого продукта {[}18{]}.

В процессе биосинтеза новый продуцент F-116 образует бактериоцин,
отличным от известного бактериоцина низина. Штамм генотипирован по гену
16S pPНК (\emph{GenBank} EF 100777). На основе штамма F-116 создан
препарат LGS, широкого спектра антимикробного действия, эффективный
против патогенных и условно-патогенных контаминантов, развивающихся в
продуктах питания и сырье при хранении. Активно формируется новое
научно-практическое направление -- разработка рецептур консервирующих
покрытий из нескольких биоконсервантов, обладающих полезными для
организма человека свойствами, и сочетание обработки консервантами с
низкотемпературным режимом хранения.

Таким образом, проведенный анализ литературы подчёркивает актуальность
разработки и внедрения интегрированных решений: пищевых покрытий в
сочетании с оптимальными режимами охлаждения, способных существенно
повысить сохранность мяса птицы. При этом современная научная база
нуждается в расширении сведений о механизмах взаимодействия таких
покрытий с микробиологической средой при близкокриоскопическом хранении,
что составляет основу настоящего исследования.

Целью работы является определение влияния раствора антимикробного
защитного покрытия на основе молочной сыворотки и культуральной жидкости
на потери массы тушек птицы и изменение микробиологических показателей
тушек птицы в процессе охлаждения и хранения при близкриоскопической
температуре в течение 30 суток.

Новизна исследований заключается в определении антимикробных и
влагобарьерных свойств, разработанных пищевых пленкообразующих составов
на основе природных компонентов и бактериоцина рекомбинантного штамма
лактококка.

{\bfseries Материалы и методы.} Объектом исследования послужили тушки и
полу-тушки цыплят-бройлеров, пищевые покрытия содержащие органические
кислоты, а также бактериоцин гибридного штамма \emph{Lactococcus lactis}
ssp. \emph{lactis} F-116.

На первом этапе исследований были проведены модельные эксперименты по
обоснованию оптимальной концентрации пищевого пленкообразующего состава.
Опыты проводились на серийно выпускае­мом пищевом пленкообразующем
составе (ППС-10) на основе моноглицеридов (Таблица 1). Водные растворы
ППС готовили по рецептуре с массовой долей сухих веществ 0,1; 1; 3; 5 и
10\%. Готовые растворы наносили методом распыления на модели-аналоги
мяса птицы по теплофизиче­ским свойствам, выполненных из геля агар-агара.
Экспериментальные исследования проводили на технологическом стенде
лаборатории холодильной технологии ВНИХИ - филиал ФГБНУ «ФНЦ пищевых
систем им. В.М. Горбатова» РАН.
\end{multicols}

\tcap{Таблица 1 - Рецептура пищевого пленкообразующего состава (ППС-10) на основе моноглицеридов}
\begin{longtblr}[
  label = none,
  entry = none,
  note{} = {\textit{Примечание: * - массовая доля сухих веществ приведена в пересчете на 3\% раствор ППС-10}},
]{
  cells = {font = \small},
  row{1} = {c,font=\bfseries\small},
  column{2} = {c},
  row{6} = {c,font=\itshape\small},
  cell{6}{1} = {c=2}{},
  hlines,
  vlines,
}
Наименование сырья & Кол-во, \%\\
Пищевой пленкообразующийсостав ППС-10 & 3,0\\
{\textit{Включающий:}\\\labelitemi\hspace{\dimexpr\labelsep+0.5\tabcolsep}Моноглицериды дистиллированные\\\labelitemi\hspace{\dimexpr\labelsep+0.5\tabcolsep}Эмульгатор вспомогательный по СТП Нижегородского МЖК 020500-59-86\\\labelitemi\hspace{\dimexpr\labelsep+0.5\tabcolsep}Калий сорбиновокислый или сорбиновая кислота\\\labelitemi\hspace{\dimexpr\labelsep+0.5\tabcolsep}Ацетилированные моноглицериды дистиллированные\\\labelitemi\hspace{\dimexpr\labelsep+0.5\tabcolsep}Глицерин дистиллированный} & {\\0,84*\\0,48*\\0,06*\\1,32*\\0,3*}\\
Молочная кислота & 2,0\\
Вода питьевая & 95,0\\
\end{longtblr}

\begin{multicols}{2}
На втором этапе были проведены исследования по определению влияния
пищевого покрытия на основе молочной сыворотки на потери массы птицы в
условиях холодильного хранения и на микробиологические показатели. Для
приготовления ППС на основе молочной сыворотки нагревали питьевую воду в
котле с подогревом и мешалкой до температуры 55±2 °С и при постоянном
перемешивании поочередно вносили молочную сыворотку, глицерин, молочную
кислоту, лимонную кислоту и бактериоцин гибридного штамма
\emph{Lactococcus lactis} ssp. \emph{lactis} F-116 в количестве,
указанном в таблице 2. Перемешивание продолжали до получения однородной
эмульсии. Нанесение покрытия на полу-тушки птицы осуществляли перед
охлаждением методом распыления с помощью распылителя Gigant GS-01
(температура раствора при распылении 55±2 °С).
\end{multicols}

\tcap{Таблица 2 - Рецептура пищевого пленкообразующего состава на основе молочной сыворотки}
\begin{longtblr}[
  label = none,
  entry = none,
]{
  cells = {font = \small},
  row{1} = {font = \bfseries\small},
  row{2} = {font = \small},
  column{2} = {c},
  hlines,
  vlines,
}
Наименование компонента & Кол-во, \%\\
Молочная сыворотка & 5,0\\
Молочная кислота & 2,0\\
Лимонная кислота & 0,5\\
Глицерин & 10,0\\
Культуральная жидкость \textit{Lactococcus}\textit{ }\textit{lactis}\textit{ }ssp.\textit{ }\textit{lactis} F-116 & 23,0*\\
Вода & 80,2
\end{longtblr}

\emph{*Содержание сухих веществ в культуральной жидкости 10,0 \%.}

\begin{multicols}{2}
Контрольные образцы полу-тушек птицы исследовали после охлаждения и
через 5, 10, 15 суток хранения в упакованном виде, а опытные образцы
исследовали после охлаждения, упаковки и через 5, 10, 15, 20, 25, 30
суток хранения. Образцы хранили при температуре 0±1 °С, относительной
влажности 8590\%, до появления явных признаков порчи по внешнему виду
образцов, органолептическим и санитарно-гигиеническим показателям.
Потери массы определяли взвешиванием образцов на весах CAS SW-5.

Отбор проб и подготовку к микробиологическим исследованиям проводили
согласно ГОСТ 31467-2012 (не менее 3-х проб). Общее микробное число
(ОМЧ) определяли по ГОСТ Р 7702.2.1-2017 (Продукты убоя птицы, продукция
из мяса птицы и объекты окружающей производственной среды).

На третьем этапе наших исследований были проведены эксперименты по
определению влияния антимикробной активности коммерческого препарата
«Nisin Рro» в сочетании с культуральной жидкостью штамма
\emph{Lactococcus lactis} ssp. \emph{lactis} F-116 на общее микробное
число образцов мяса птицы в условиях холодильного хранения. В качестве
основного экспериментального образца был приготовлен раствор ППС,
состоящий из: 0,1\%, культуральной жидкости (КЖ) штамма F-116, 0,1\%
раствора низина, 1\% желатина и 2\% молочной кислоты.

Все опыты имели не менее 3-х повторностей, обеспечивающих получение
воспроизводимых и достоверных данных. Результаты, полученные в ходе
эксперимента, с оформлением таблиц, обрабатывали с помощью компьютерных
программ Excel (Microsoft Inc.) и Statistica for Windows (v.5.0
StatSoftInc), рассчитывая среднее арифметическое, доверительные
интервалы, стандартное отклонение. Достоверность различий между средними
величинами оценивали с использованием t-критерия Стьюдента (р≤0,05).

{\bfseries Результаты и обсуждение.} Экспериментальные данные по влиянию
различных концентраций ППС-10 на уменьшение усушки и изменение
коэффициента сопротивления испа­рению влаги при охлаждении и по­следующем
хранении проводили в модельных опытах в течение 5 суток (рис.1). Из
данных рисунка следует, что величина усушки опытных образцов
уменьшается.
\end{multicols}

\fig[0.6\textwidth]{p2/image22}[Рис.1 - Изменение коэффициента
сопротивления испарению воды с поверхности моделей в зависимости от
концентрации раствора ППС для продолжительности их хранения 120 часов]

\begin{multicols}{2}
С ростом концентрации от 0 до 10\%, а вели­чина сопротивления испарению
(R) влаги увеличивается с рос­том концентрации раствора пищевого покрытия
и достигает значения R≈2 при концентрации ППС в растворе 3-5\%. Разница
между значениями потерь массы опытных образцов для концентрации 3 и 5\%
недостоверна при p≤0,05. Это дает возможность сделать вывод о том, что
для ППС-10 на основе моноглицеридов с концентрацией выше 3\% применять
экономически нецелесообразно.

В предыдущих исследованиях {[}19{]} ППС-10 в сочетании с молочной
кислотой показал весьма высокий уровень антимикробной защиты, однако
недостатком данного состава является сложность его получения, в связи с
чем был создан новый пищевой пленкообразующий состав на основе природных
органических компонентов.

Экспериментальный пищевой пленкообразующий состав (ППС) приготовлен на
основе сухой молочной сыворотки (ГОСТ 33958-2016), которая содержит
биологически активные вещества, включая белки, углеводы, незаменимые
аминокислоты, витамины А, С, Е, Н и группы В, биотин, холин,
микроэлементы-калий, кальций, фосфор, магний, йод и цинк и молочнокислые
бактерии (до ≤1×10² КОЕ/г). Лимонная кислота играет роль антиокислителя,
предотвращая окисление подкожных жиров, содержащихся в тушках птицы.
Глицерин служит в пищевом пленкообразующем составе пластификатором.

В качестве природного нетоксичного антимикробного агента использовали
препарат LGS, активным веществом которого служит бактериоцин по
типу низина, синтезируемый рекомбинантным штаммом \emph{L. lactis} ssp.
\emph{lactis} F-116 (GenBank: EF100777.1) при культивировании в
ферментационной среде, следующего состава (г/л): сахароза - 20,0;
\(KH_{2}PO_{4}\)- 20,0; NaCl - 2,0; \(MgSO_{4}\)- 0,2; дрожжевой
экстракт - 20,0; рН 6,8 - 7,0. Лиофильно высушенный препарат обладает
активностью 1100 МЕ/мг в пересчете на активность эталонного препарата
``Nisaplin'' фирмы`` Aplin \& Barrett, LTD.

Пик бактериоцинсинтезирующей активности штамма F-116 приходится на 12-й
час культивирования. Важно учитывать особенности локализации целевого
продукта: в то время как часть бактериоцина экстрагируется в среду
непосредственно в процессе ферментации, до 50\% его объема остается
прочно связанным с мембранами и клеточными стенками продуцента.

Результаты экспериментальных исследований по определению влияния
обработки поверхности тушек птицы антимикробным защитным покрытием на
основе молочной сыворотки и культуральной жидкости штамма F-116 на
товарный вид и потери массы в процессе охлаждения и хранения при
близкриоскопической температуре представлены на рисунках 2 и 3.
\end{multicols}

\begin{figs}[Рис.2 - Товарный вид охлажденного мяса птицы после
обработки: \\а - до упаковки; б - после упаковки.]
\fig[0.47\textwidth][5.5cm]{p2/image13}[а)]
\fig[0.47\textwidth][5.5cm]{p2/image14}[б)]
\end{figs}

\fig[0.7\textwidth]{p2/image23}[Рис.3 - Влияние охлаждения и хранения
на потери массы контрольных (1) и опытных (2) полу-тушек цыплят
бройлеров]

\begin{multicols}{2}
Графические данные (рис.3) отчетливо демонстрируют существенный разрыв
в динамике потери массы между контрольными и опытными образцами уже на
этапе охлаждения, что обусловлено массой нанесенного покрытия и его
защитными свойствами, препятствующими испарению влаги с поверхности
птицы.

Важно отметить: преимущество опытных образцов со временем только
усиливается. Если сразу после охлаждения разница в усушке составляла
0,45\%, то к 10-м суткам она достигла 0,60\%, увеличившись в 1,3 раза, в
чем и заключается принципиальное отличие предлагаемой технологии от
традиционных методов интенсивного охлаждения (воздушного или
гидроаэрозольного). При использовании стандартных способов разница в
массе между контролем и опытом в процессе хранения обычно сглаживается,
тогда как применение ППС позволяет стабильно удерживать влагу. В
конечном итоге, показатели массы опытных полутушек после 15 суток
хранения оказались даже на 0,20\% выше, чем у контрольных образцов,
хранившихся всего 10 суток.

Полученные данные подтверждают, что основной разрыв в показателях убыли
массы между контрольными и опытными образцами формируется уже на этапе
охлаждения, что может быть обусловлено двумя факторами: во-первых,
изменением исходного веса птицы за счет нанесения покрытия, и во-вторых
--- созданием физического барьера, который значительно повышает
коэффициент сопротивления испарению влаги. Важно подчеркнуть, что при
использовании ППС разница в усушке между контролем и опытом со временем
только нарастает. Если сразу после охлаждения этот показатель составляет
0,45\%, то к 15-м суткам он достигает 0,65\%, увеличиваясь в 1,4 раза. В
этом заключается принципиальное отличие предлагаемого метода от
интенсивных способов охлаждения (при низких температурах и высокой
скорости воздуха) или гидроаэрозольной обработки: в тех случаях разница
в массе между образцами в процессе последующего хранения, напротив,
имеет тенденцию к сокращению.

В наших экспериментах общее микробное число (ОМЧ) в мышечной ткани
полу-тущек цыплят контрольных образцов (образцы обработаны водой)
составило 2,6±0,08×10\tsp{3} КОЕ/г, а обработанных ППС +
бактериоцин 3,2±0,01×10\tsp{3}, продукт по показателям
качества соответствовал требования СанПиН 2.3.2.1078-01 (Гигиенические
требования безопасности и пищевой ценности пищевых продуктов (Табл.3).
\end{multicols}

\tcap{Таблица 3 - Изменение общего микробного числа (ОМЧ) после обработки тушек птицы консервирующими агентами и хранения в условиях холодильной камеры при 0±1℃}
\begin{longtblr}[
  label = none,
  entry = none,
  note{} = {\emph{Примечание: ППС - пищевой пленкообразующий состав на основе молочной сыворотки}},
]{
  width = \linewidth,
  colspec = {Q[20]Q[350]Q[131]Q[121]Q[152]Q[117]},
  cells = {c, font = \small},
  row{1} = {font=\bfseries\small},
  row{2} = {font=\bfseries\small},
  cell{1}{1} = {r=2}{},
  cell{1}{2} = {r=2}{},
  cell{1}{3} = {c=4}{0.521\linewidth},
  hlines,
  vlines,
}
№ & Cостав консервирующей композиции & Общее микробное число после обработки и хранения, КОЕ/г &  &  & \\
 &  & исходное & 10 суток & 15 суток & 20 суток\\
1 & ППС+0,1\% бактериоцин ЛГС + 2\% молочной кислоты & 3,2±0,01×10\textsuperscript{3} & 4,7±0,02×10\textsuperscript{3} & 2,4 ±0,09×10\textsuperscript{4} & 8,7±0,05×10\textsuperscript{4}\\
2 & Контроль (вода питьевая) & 2,6±0,08×10\textsuperscript{3} & 7,5±0,04×10\textsuperscript{5} & Сняты с хранения & -
\end{longtblr}

\begin{multicols}{2}
Следует отметить, что микробная обсемененность образцов, обработанных
ППС на основе молочной сыворотки и с препаратом LGS, на первом этапе
хранения одинаков и представлен, в основном, мелкими колониями,
характерными для молочнокислых бактерий. Эффект обусловлен наличием в
препарате LGS лактококков --- активных продуцентов бактериоцинов. Данные
микроорганизмы способны восстанавливаться после стадии анабиоза, в
которую они переходят при лиофильной сушке. Однако при длительном
воздействии низких температур популяция мезофильных лактококков
постепенно сокращается, а их место занимают психротрофные
микроорганизмы, обладающие природной устойчивостью к действию низина.

В процессе холодильного хранения состав микробиоты, колонизирующей
поверхность тушек, закономерно трансформируется. На фоне отмирания
вегетативных мезофильных форм (включая и сами лактококки) наблюдается
ускоренное развитие психротрофной флоры. Тем не менее использование
защитного состава позволило существенно замедлить темпы микробной
обсемененности. Согласно данным таблицы 2, через 10 суток хранения
показатель ОМЧ в опытных образцах увеличился всего в 1,5 раза, в то
время как в контрольной группе рост был практически трехсоткратным.

В контрольном образце после 10 суток хранения в условиях охлаждения ОМЧ
за 10 суток хранения составило 7,5±0,04×10\tsp{5}. Эти
образцы мяса птицы (обработанные водой) на 10 сутки имели
органолептические показатели, присущие испорченному мясу птицы
(ослизнение поверхности, неприятный запах{\bfseries ).} В образцах мяса
были обнаружены бактерии из рода \emph{Pseudomonas,} энтерококки,
анаэробные гнилостные бактерии \emph{Proteus s}pp., вызывающие
неприятный запах, ослизнение поверхности образцов, были обнаружены
дрожжи и плесени, показатели качества мяса птицы превосходили допустимый
уровень микробиальной и органолептической безопасности по СанПиН
2.3.2.1078-01. Образцы сняты с хранения.

В опытных образцах мяса птицы, обработанных пищевым пленкообразующим
составом, общее микробное число на 10 сутки хранения составило
4,7±0,02×10\tsp{3} КОЕ/г, на 15 сутки
2,4±0,09×10\tsp{4} КОЕ/г, а на 20 сутки
8,7±0,05×10\tsp{4} КОЕ/г, что не превышало предельно
допустимое значение микробиологической обсемененности в соответствии с
требованием СанПиН 2.3.2.1078-01.

Результаты исследований по изучению влияния обработки тушек птицы
антимикробным защитным покрытием на основе 1\% желатина, 0,1\% препарата
``Nisin Рro'' и культуральной жидкости рекомбинантного штамма F-116 в
процессе хранения при близкриоскопической температуре на общее микробное
число (ОМЧ) показали, что по мере хранения в контрольном образце,
обработанном питьевой водой, без добавления консервирующих агентов,
наблюдали постепенное увеличение обсемененности: исходное количество
микробов в этом образце от 2,6×10\tsp{3} КОЕ/г за 5 дней
увеличилось на порядок, а за 10 суток до 3,9×10\tsp{5}
КОЕ/г, что превышает допустимый уровень обсемененности охлажденной
птицы, а за последующие 15 дней хранения при 0±1 °С ОМЧ достигло
1,2×10\tsp{8} КОЕ/г. Отмечено ослизнение поверхности
образцов мяса, цвет сероватый, затхлость и неприятный запах, были
обнаружены бактерии бактерии \emph{Proteus vulgaris}, плесени и дрожжи,
являющиеся преобладающими контаминантами, ответственных за порчу мяса
птицы, которые преимущественно находятся на его поверхности,~а,
размножаясь в сырье, вызывают биохимические реакции, которые приводят к
таким изменениям, как протеолиз, деградация аминокислот и самое главное,
окисление липидов, которые на поздних стадиях хранения мяса птицы
вызывают появление посторонних привкусов и слизи, что приводят к
снижению качества и сокращению коммерческого срока годности продукта.

В образцах, обработанных 0,1\% раствором «Nisin Рro», уже на 10-е сутки
хранения показатель ОМЧ превысил предельно допустимые значения (Табл.
4).
\end{multicols}

\tcap{Таблица 4 - Изменение общего микробного числа (ОМЧ) после обработки тушек цыплят препаратом ``Nisin Рro'' в сочетании с желатином, молочной кислотой и культуральной жидкостью штамма Lactococcus lactis ssp. lactis F-116 в процессе хранения при близкриоскопической температуре в течение 30 суток}
\begin{longtblr}[
  label = none,
  entry = none,
]{
  width = \linewidth,
  colspec = {Q[10]Q[170]Q[100]Q[90]Q[100]Q[100]Q[100]Q[100]Q[100]},
  cells = {c, font = \scriptsize},
  row{1} = {font=\bfseries\scriptsize},
  row{2} = {font=\bfseries\scriptsize},
  cell{1}{1} = {r=2}{},
  cell{1}{2} = {r=2}{},
  cell{1}{3} = {c=7}{0.695\linewidth},
  hlines,
  vlines,
}
№ & Добавки к составу ППС & Общее микробное число после обработки и хранения (КОЕ/г) &  &  &  &  &  & \\
 &  & Начальное (после обработки) & 5 сут. & 10 сут. & 15сут. & 20 сут. & 25 сут. & 30 сут.\\
1 & 0,1\% “Nisin Рro” & 2,3±0,01×10\tsp{3} & 7,2±0,03×10\tsp{3} & 2,7±0,01×10\tsp{4} & 6,5±0,03×10\tsp{5} & 2,1±0,03×10\tsp{6} & 4,8±0,04×10\tsp{6} & 7,1±0,03×10\tsp{7}\\
2 & 0,1\%“Nisin Рro” + 1-\% желатина & 3,1±0,02×10\tsp{3} & 5,4±0,03×10\tsp{3} & 7,3±0,03×10\tsp{4} & 5,3±0,03×10\tsp{5} & 1,2±0,03×10\tsp{6} & 4,8±0,04×10\tsp{6} & 4,1±0,02×10\tsp{7}\\
3 & “Nisin Рro” + 2\% МК & 2,5±0,01×10\tsp{3} & 4,6±0,02×10\tsp{3} & 2,8±0,02×10\tsp{4} & 8,6±0,05×10\tsp{4} & 1,4±0,03×10\tsp{4} & 1,4±0,02×10\tsp{6} & 3,7±0,01×10\tsp{7}\\
4 & 0,1\%“NisinРro” +1\% желатина + 2\% МК & 3,0±0,02×10\tsp{3} & 4,2±0,02×10\tsp{3} & 5,1±0,02×10\tsp{3} & 2,7±0,02×10\tsp{4} & 4,1±0,02×10\tsp{5} & 1,6±0,01×10\tsp{6} & 8,6±0,04×10\tsp{6}\\
5 & 0,1\% “Nisin Рro” +1\% желатина + 2\% МК + 0,1\% КЖ \textit{L.lactis} & 5,4±0,02×10\tsp{3} & 9,1±0,09×10\tsp{3} & 1,4±0,003×10\tsp{4} & 5,1±0,03×10\tsp{4} & 9,4±0,04×10\tsp{4} & 1,1±0,02×10\tsp{5} & 1,1±0,02×10\tsp{6}\\
6 & Контроль (вода питьевая) & 2,6±0,01×10\tsp{3} & 2,7±0,01×10\tsp{4} & 3,9±0,02×10\tsp{5} & 7,0±0,03×10\tsp{6} & 8,6±0,04×10\tsp{6} & 5,7±0,03×10\tsp{7} & 1,2±0,05×10\tsp{8}
\end{longtblr}

\begin{multicols}{2}
При анализе комбинированного воздействия (0,1\% низина и 1\% желатина)
было отмечено, что на начальном этапе уровень обсемененности тушек
оказался несколько выше по сравнению с контролем, обработанным питьевой
водой. Вероятно, это связано с исходным микробиологическим фоном самого
желатина. Тем не менее в процессе дальнейшего хранения динамика роста
микроорганизмов заметно замедлилась. Наблюдаемое снижение численности
мезофильных бактерий подтверждает эффективность желатина в качестве
защитного барьера, препятствующего повторной контаминации поверхности
продукта.

Общая обсемененность образцов, обработанных 0,1\%-ным раствором
препарата ``Nisin Рro'' с добавлением 2\% молочной кислоты, увеличивает
срок хранения образцов до 15 суток при сохранении органолептических
показателей. Сочетание действия низина, желатина, молочной кислоты и
культуральной жидкости \emph{L. lactis ssp.} \emph{lactis} F-116 -
продуцента бактериоцина широкого спектра действия позволило увеличить
срок хранения до 20 суток в процессе хранения при близкриоскопической
температуре, дрожжевые клетки не были обнаружены, и по внешнему виду и
органолептическим показателям мясо соответствовало ГОСТ 10444.13.
Исследования показали, что ни в одном из образцов патогенные
микроорганизмы рода \emph{Salmonella}, а также \emph{Listeria
monocytogenes} не были обнаружены.

{\bfseries Выводы.} Установлена зависимость между концентрацией
пленкообразующего состава на основе моноглицеридов (ППС-10) и
коэффициентом сопротивления испарению влаги. На основе полученных данных
обоснована оптимальная концентрация компонентов, обеспечивающая
максимальную защиту поверхности птицы в процессе охлаждения и
последующего хранения.

Установлено, что применение ППС на базе молочной сыворотки и
культуральной жидкости позволяет существенно минимизировать усушку. Если
сразу после охлаждения разница в потерях массы между опытом и контролем
составляла 0,45\%, то к 15-м суткам хранения разрыв увеличился в 1,4
раза (до 0,65\%), что доказывает преимущество ППС перед интенсивными
методами охлаждения (воздушным или гидроаэрозольным), при которых
разница в массе между образцами в процессе хранения обычно нивелируется.

Экспериментально подтверждено, что использование защитных покрытий
обеспечивает удержание влаги в тканях. Потери массы опытных образцов
через 15 суток хранения в упаковке оказались на 0,20\% ниже, чем у
контрольных тушек на 10-е сутки. В то же время традиционные способы
высокоскоростного охлаждения не обеспечивают столь длительного
барьерного эффекта.

Доказано, что обработка птицы многокомпонентным составом (молочная
сыворотка, глицерин, органические кислоты и бактериоцин штамма
\emph{Lactococcus lactis} ssp. \emph{lactis} F-116) эффективно подавляет
развитие контаминантной микробиоты. Внедрение данной технологии
позволяет увеличить срок хранения продукции при температуре 0±1 °C до 20
суток.

Использование натуральных биоконсервантов не только гарантирует
микробиологическую безопасность, но и повышает биологическую ценность
продукта. Обогащение мяса аминокислотами, витаминами и антиоксидантами,
входящими в состав сыворотки и бактериоцинов, открывает перспективы для
производства функциональных продуктов питания, способствующих укреплению
здоровья потребителей.

Установлено, что сформированный на поверхности барьер препятствует
проникновению патогенов в глубокие слои мышечной ткани и ингибирует их
размножение. Это позволяет снизить общую обсемененность (ОМЧ) на три
порядка, обеспечивая полное соответствие продукции требованиям СанПиН
2.3.2.1078-01 в течение всего периода хранения.

Сравнительный анализ показал, что комбинация низина с желатином и 2\%
молочной кислотой обладает умеренным защитным действием, но недостаточно
эффективно сдерживает рост дрожжей и гнилостных бактерий при
низкотемпературных режимах. Полученные результаты указывают на
выраженный синергетический потенциал при совместном использовании низина
с другими антимикробными агентами и бактериоцинами различного спектра
действия.

Результаты апробации в условиях реального производства подтвердили
высокую эффективность разработанной технологии. Её внедрение позволяет
одновременно решить три задачи: повысить санитарно-гигиеническую
безопасность продукции, увеличить сроки её реализации и существенно
сократить экономические потери, связанные с естественной убылью массы.
\end{multicols}

\begin{center}
{\bfseries Литература}
\end{center}

\begin{refs}
1. Fabra M. J., Talens P., Chiralt A. Influence of calcium on tensile,
optical and water vapour permeability properties of sodium caseinate
edible films //Journal of Food Engineering. - 2010.-Vol.96(3). - P.
356-364. DOI 10.1016/j.jfoodeng.2009.08.010.

2. Kumar L. et al. Edible films and coatings for food packaging
applications: A review //Environmental Chemistry Letters.
-2021.-Vol.20(4). -P.875-900. DOI 10.1007/s10311-021-01339-z.

3. Ramos O. L. et al. Edible films and coatings from whey proteins: a
review on formulation, and on mechanical and bioactive properties
//Critical reviews in food science and nutrition. - 2012.-Vol.52(6).
-P.533-552. DOI 10.1080/10408398.2010.500528.

4. Elsabee M. Z., Abdou E. S. Chitosan based edible films and coatings:
A review //Materials science and engineering: С.-2013. Vol.33(4).
-P.1819-1841.DOI 10.1016/j.msec.2013.01.010.

5. Sánchez-Ortega I., Blanca E. Garcia-Almendarez, et al.
Characterization and antimicrobial effect of starch-based edible coating
suspensions//Food Hydrocolloids. -2016.- Vol.52.-P.906-913. DOI
10.1016/j.foodhyd.2015.09.004.

6. Li X. Zhand R, et al. Active packaging for the extended shelf-life of
meat: perspectives from consumption habits, market requirements and
packaging practices in China and New Zealand //Foods. -2022.-Vol.11(18):
2903. DOI 10.3390/foods11182903.

7. Yi J, Chen F, Xu C.et al. Decoding spoilage metabolism in chicken
breast: A functional microbial perspective on off-flavor
formation//International Journal of Food Microbiology. -
2025. -Vol.442:111367. DOI 10.1016/j.ijfoodmicro.2025.111367.

8. Huezo R. Smith D.P.et al. Effect of immersion or dry air chilling on
broiler carcass moisture retention and breast fillet functionality
//Journal of applied poultry research. - 2007.-Vol.16(3). -P.438-447.
DOI 10.1093/japr/16.3.438.

9. James C., Vincent C. et al. The primary chilling of poultry
carcasses-a review Refrigeration des carasses de poulet:
tendances//International Journal of Refrigeration. - 2006. -Vol.29(6).
-P.847-862. DOI 10.1016/j.ijrefrig.2005.08.003.

10. Jeong J. Y. Janardhan K.K. et al. Moisture content, processing
yield, and surface color of broiler carcasses chilled by water, air, or
evaporative air//Poultry science. -2011.-Vol.90(3). -P.687-693. DOI
10.3382/ps.2010-00980.

11. Рогов И.А. Технология мяса и мясных продуктов. //КолосС. -2009.
--С.711. ISBN 978-5-9532-0644-0.

12. В. М. Позняковский, О. А. Гончаренко, Д. В. Криштафович.
Товароведение и экспертиза мясных и мясосодержащих продуктов.
//Издательство "Лань». - 2020. -№ 4. -С.432.ISBN 978-5-8114-4942-2.

13. Han J. H. Chapter 9-Edible films and coatings: a review //Innovations
in food packaging (Second Edition). -2014. -P.213-255. DOI
10.1016/B978-0-12-394601-0.00009-6.

14. Erdem Büyükkiraz M., Kesmen Z. Antimicrobial peptides (AMPs): A
promising class of antimicrobial compounds//Journal of applied
microbiology. -2022. -Vol.132(3). -P.1573-1596. DOI 10.1111/jam.15314.

15. Sánchez-Ortega I. Et al. Antimicrobial edible films and coatings for
meat and meat products preservation//The scientific world journal.
-2014. -Vol.2014(1)/248935. DOI 10.1155/2014/248935.

16. Khare A. K., Babu R.N, et al. Utilization of carrageenan, citric acid
and cinnamon oil as an edible coating of chicken fillets to prolong its
shelf life under refrigeration conditions //Veterinary World. -2016.
-Vol.9(2):166-75. DOI 10.14202/vetworld.2016.166-175.

17. Moura-Alves M. Et al. Antimicrobial and antioxidant edible films and
coatings in the shelf-life improvement of chicken meat//Foods.
-2023.-Vol.12(12):2308. DOI 10.3390/foods12122308.

18. Stoyanova L. G., Napalkova M. V., Netrusov A. I. The creating a new
biopreservative based on fusant strain Lactococcus lactis SSP. lactis
F-116 for food quality and its safety //Hygienic Engineering and Design.
-2016.-Vol.16.-P.19-27. -URL:
\url{https://www.cabidigitallibrary.org/doi/full/10.5555/20173079592.-}
Date obrashcheniya: 07.03.2026.

19. Umiraliyeva L., Stoyanova L. Dibirasulaev M, et al. Study of the
effect of food coatings based on monoglycerides with the inclusion of
biopreservatives in their composition on the microbiological safety
during storage of chilled poultry meat //CyTA-Journal of Food.
-2025.-Vol.23(1):2561649. DOI 10.1080/19476337.2025.2561649.
\end{refs}

\begin{center}
{\bfseries References}
\end{center}

\begin{refs}
1. Fabra M. J., Talens P., Chiralt A. Influence of calcium on tensile,
optical and water vapour permeability properties of sodium caseinate
edible films //Journal of Food Engineering. - 2010.-Vol.96(3). - P.
356-364. DOI 10.1016/j.jfoodeng.2009.08.010.

2. Kumar L. et al. Edible films and coatings for food packaging
applications: A review //Environmental Chemistry Letters.
-2021.-Vol.20(4). -P.875-900. DOI 10.1007/s10311-021-01339-z.

3. Ramos O. L. et al. Edible films and coatings from whey proteins: a
review on formulation, and on mechanical and bioactive properties
//Critical reviews in food science and nutrition. - 2012.-Vol.52(6).
-P.533-552. DOI 10.1080/10408398.2010.500528.

4. Elsabee M. Z., Abdou E. S. Chitosan based edible films and coatings:
A review //Materials science and engineering: С.-2013. Vol.33(4).
-P.1819-1841.DOI 10.1016/j.msec.2013.01.010.

5. Sánchez-Ortega I., Blanca E. Garcia-Almendarez, et al.
Characterization and antimicrobial effect of starch-based edible coating
suspensions//Food Hydrocolloids. -2016.- Vol.52.-P.906-913. DOI
10.1016/j.foodhyd.2015.09.004.

6. Li X. Zhand R, et al. Active packaging for the extended shelf-life of
meat: perspectives from consumption habits, market requirements and
packaging practices in China and New Zealand //Foods. -2022.-Vol.11(18):
2903. DOI 10.3390/foods11182903.

7. Yi J, Chen F, Xu C.et al. Decoding spoilage metabolism in chicken
breast: A functional microbial perspective on off-flavor
formation//International Journal of Food Microbiology. -
2025. -Vol.442:111367. DOI 10.1016/j.ijfoodmicro.2025.111367.

8. Huezo R. Smith D.P.et al. Effect of immersion or dry air chilling on
broiler carcass moisture retention and breast fillet functionality
//Journal of applied poultry research. - 2007.-Vol.16(3). -P.438-447.
DOI 10.1093/japr/16.3.438.

9. James C., Vincent C. et al. The primary chilling of poultry
carcasses-a review Refrigeration des carasses de poulet:
tendances//International Journal of Refrigeration. - 2006. -Vol.29(6).
-P.847-862. DOI 10.1016/j.ijrefrig.2005.08.003.

10. Jeong J. Y. Janardhan K.K. et al. Moisture content, processing
yield, and surface color of broiler carcasses chilled by water, air, or
evaporative air//Poultry science. -2011.-Vol.90(3). -P.687-693. DOI
10.3382/ps.2010-00980.

11. Rogov I.A. Tehnologija mjasa i mjasnyh produktov. //KolosS. -2009.
--S.711. ISBN 978-5-9532-0644-0. {[}in Russian{]}

12. V. M. Poznjakovskij, O. A. Goncharenko, D. V. Krishtafovich.
Tovarovedenie i jekspertiza mjasnyh i mjasosoderzhashhih produktov.
//Izdatel' stvo "Lan'». - 2020. -№ 4.
-S.432.ISBN 978-5-8114-4942-2. {[}in Russian{]}

13. Han J. H. Chapter 9-Edible films and coatings: a review //Innovations
in food packaging (Second Edition). -2014. -P.213-255. DOI
10.1016/B978-0-12-394601-0.00009-6.

14. Erdem Büyükkiraz M., Kesmen Z. Antimicrobial peptides (AMPs): A
promising class of antimicrobial compounds//Journal of applied
microbiology. -2022. -Vol.132(3). -P.1573-1596. DOI 10.1111/jam.15314.

15. Sánchez-Ortega I. Et al. Antimicrobial edible films and coatings for
meat and meat products preservation//The scientific world journal.
-2014. -Vol.2014(1)/248935. DOI 10.1155/2014/248935.

16. Khare A. K., Babu R.N, et al. Utilization of carrageenan, citric acid
and cinnamon oil as an edible coating of chicken fillets to prolong its
shelf life under refrigeration conditions //Veterinary World. -2016.
-Vol.9(2):166-75. DOI 10.14202/vetworld.2016.166-175.

17. Moura-Alves M. Et al. Antimicrobial and antioxidant edible films and
coatings in the shelf-life improvement of chicken meat//Foods.
-2023.-Vol.12(12):2308. DOI 10.3390/foods12122308.

18. Stoyanova L. G., Napalkova M. V., Netrusov A. I. The creating a new
biopreservative based on fusant strain Lactococcus lactis SSP. lactis
F-116 for food quality and its safety //Hygienic Engineering and Design.
-2016.-Vol.16.-P.19-27. -URL:
\url{https://www.cabidigitallibrary.org/doi/full/10.5555/20173079592} .-
Date obrashcheniya: 07.03.2026.

19. Umiraliyeva L., Stoyanova L. Dibirasulaev M, et al. Study of the
effect of food coatings based on monoglycerides with the inclusion of
biopreservatives in their composition on the microbiological safety
during storage of chilled poultry meat //CyTA-Journal of Food.
-2025.-Vol.23(1):2561649. DOI 10.1080/19476337.2025.2561649.
\end{refs}

\begin{info}
\hspace{1em}\emph{{\bfseries Сведения об авторах}}

Умиралиева Л.Б. - кандидат технических наук, ассоциированный
профессор, ТОО «Казахский научно-исследовательский институт
перерабатывающей и пищевой промышленности», Алматы, Казахстан, e-mail:
l.umiraliyeva@rpf.kz;

Дибирасулаев М.А.- доктор технических наук, ВНИХИ - филиал ФГБНУ «ФНЦ
пищевых систем им. В. М. Горбатова» РАН, Москва, Россия, e-mail:
dmama42@mail.ru;

Дибирасулаев Д.М.- кандидат технических наук, ВНИХИ - филиал ФГБНУ
«ФНЦ пищевых систем им. В. М. Горбатова» РАН, Москва, Россия, e-mail:
dibdm@mail.ru;

Стоянова Л.Г.- -доктор биологических наук, профессор, МГУ им. М. В.
Ломоносова г. Москва, Россия, e-mail: stoyanovamsu@mail.ru;

Турлыбаева С.С.-кандидат сельскохозяйственных наук, ТОО «Казахский
научно-исследовательский институт перерабатывающей и пищевой
промышленности», Алматы, Казахстан, e-mail: s.turlybaeva@rpf.kz;

Филатов И.Д PhD-докторант, ТОО «Казахский научно-исследовательский
институт перерабатывающей и пищевой промышленности» Алматы, Казахстан,
e-mail: filatov-vanya@inbox.ru.

\hspace{1em}\emph{{\bfseries Information about the authors}}

Umiraliyeva L.B. - Candidate of Technical Sciences, Associate
Professor, ``Kazakh Scientific Research Institute of Processing and
Food Industry'' LLP, Almaty, Kazakhstan, e-mail: l.umiraliyeva@rpf.kz,

Dibirasulaev M.A. - Doctor of Technical Sciences, All-Russian
Scientific Research Institute of Refrigeration Industry - branch of
V.M. Gorbatov Federal Research Center for Food Systems of Russian
Academy of Science, Moscow, Russia, e-mail: dmama42@mail.ru;

DibirasulaevD.M. - Candidate of Technical Sciences, All-Russian
Scientific Research Institute of Refrigeration Industry - branch of
V.M.  Gorbatov Federal Research Center for Food Systems of Russian
Academy of Science, Moscow, Russia, e-mail: dibdm@mail.ru;

Stoyanova L. G. - Doctor of Biological Sciences, Professor, M.V.
Lomonosov Moscow State University,, Moscow, Russia, e-mail:
stoyanovamsu@mail.ru;

Turlybaeva S.S. - Candidate of agricultural sciences, `Kazakh
Scientific Research Institute of Processing and Food Industry' LLP,
Almaty, Kazakhstan, e-mail: s.turlybaeva@rpf.kz;

Filatov I.D. - PhD student, ``Kazakh Scientific Research Institute of
Processing and Food Industry'' Almaty, Kazakhstan, LLP e-mail:
filatov-vanya@inbox.ru
\end{info}
